% Figure from MATH 590 notes, section 9. Compile with the notes preamble.
\begin{tikzpicture}
  % X
  \draw[black!55, thick] plot[smooth cycle, tension=0.8] coordinates {(-4.4,0) (-3.9,1.3) (-2.4,1.55) (-1.0,1.2) (-0.7,-0.1) (-1.3,-1.35) (-2.8,-1.5) (-4.1,-1.1)};
  \node[font=\small, black!70] at (-4.25,1.3) {$X$};
  % preimage in two pieces
  \fill[figB!12] plot[smooth cycle, tension=0.9] coordinates {(-3.9,0.35) (-3.5,1.0) (-2.8,0.9) (-2.7,0.25) (-3.3,-0.1)};
  \draw[figB, thick, dashed] plot[smooth cycle, tension=0.9] coordinates {(-3.9,0.35) (-3.5,1.0) (-2.8,0.9) (-2.7,0.25) (-3.3,-0.1)};
  \fill[figB!12] plot[smooth cycle, tension=0.9] coordinates {(-2.3,-0.7) (-1.9,-0.1) (-1.25,-0.3) (-1.3,-1.0) (-1.9,-1.15)};
  \draw[figB, thick, dashed] plot[smooth cycle, tension=0.9] coordinates {(-2.3,-0.7) (-1.9,-0.1) (-1.25,-0.3) (-1.3,-1.0) (-1.9,-1.15)};
  \node[font=\scriptsize, figB] at (-3.3,-0.75) {$f^{-1}(V)$};
  % Y
  \draw[black!55, thick] plot[smooth cycle, tension=0.8] coordinates {(1.0,0) (1.4,1.3) (2.8,1.5) (4.0,1.0) (4.3,-0.2) (3.7,-1.35) (2.2,-1.45) (1.2,-1.0)};
  \node[font=\small, black!70] at (4.2,1.25) {$Y$};
  \fill[figR!15] plot[smooth cycle, tension=0.9] coordinates {(2.0,0.1) (2.4,0.85) (3.3,0.8) (3.6,0.0) (3.1,-0.7) (2.3,-0.6)};
  \draw[figR, thick, dashed] plot[smooth cycle, tension=0.9] coordinates {(2.0,0.1) (2.4,0.85) (3.3,0.8) (3.6,0.0) (3.1,-0.7) (2.3,-0.6)};
  \node[font=\scriptsize, figR] at (3.35,-0.95) {$V$};
  % a point and its image
  \fill[black] (-3.3,0.5) circle (1.4pt) node[left=1pt, font=\scriptsize] {$x$};
  \fill[black] (2.75,0.35) circle (1.4pt) node[below=1pt, font=\scriptsize] {$f(x)$};
  \draw[-stealth, black!40, thin] (-3.2,0.6) to[bend left=18] (2.65,0.45);
  % the map
  \draw[-stealth, thick] (-0.35,1.7) to[bend left=15] node[above, font=\scriptsize] {$f$} (0.9,1.7);
\end{tikzpicture}
